\item Reconoce saberes ancestrales acerca del espacio geográfico, formas de ubicación y representaciones en México y el mundo.
\item Comprende las categorías de análisis espacial para explicar las características del espacio geográfico: lugar, región, paisaje y territorio.
\item Utiliza los conceptos de localización, distribución, diversidad, temporalidad y cambio, así como relación e interacción para el estudio del espacio geográfico.
\item Interpreta representaciones cartográficas para obtener información de diversos lugares, regiones, paisajes y territorios.
\item Emplea recursos tecnológicos para obtener y representar información geográfica en las escalas local, nacional y mundial.
\item Identifica qué son las placas tectónicas, cuáles son sus características y dinámica.
\item Argumenta la relación entre las placas tectónicas y las regiones sísmicas y volcánicas en México y el mundo para fortalecer la cultura de la prevención.
\item Relaciona los movimientos de las placas tectónicas con la distribución del relieve de la superficie terrestre y reconoce otros agentes que lo modelan.
\item Indaga sobre el origen, los usos y los problemas del suelo en la localidad.
\item Analiza la distribución de las aguas continentales en
\item México y el mundo: principales ríos, lagos, aguas subterráneas, llanuras inundables y humedales.
\item Reconoce la importancia de las cuencas hidrográficas para el desarrollo económico en México, así como para la conservación del agua y la tierra.
\item Valora el mar territorial, la zona económica exclusiva de México y sus litorales, como recursos que contribuyen al desarrollo del país, lo que fortalece la conciencia marítima.
